RESTORE reconstructs base (:term:base.terms.architectural_state:) and enabled extension state from a 4-KiB-aligned, CPUID-defined save area.

RESTORE treats its address-role EA as the base of the save area without first reading an EA-sized operand. It restores R0 through R15 and FLAGS from the fixed base block. A set \texttt{GS\_VALID} bit selects the corresponding GS image for validation and restoration; a clear bit preserves the current GSn. User-mode RESTORE ignores saved STATUS and DFA. Supervisor RESTORE validates and restores the complete saved supervisor STATUS, including EDEPTH and UO, together with current DFA as an architectural context transition. Ordinary WRSTATUS retains its defined (:term:base.terms.field:) restrictions. A set state-block bitmap bit restores its CPUID-described extension component. A clear bit for a described component installs that component\textquotesingle{}s initial state and makes the component clean; a restored non-initial image makes the component modified.

Before reading any extension component or changing (:term:base.terms.architectural_state:), RESTORE validates the complete base (:term:base.terms.header:) and bitmap. An unrecognized format, a nonzero (:term:base.terms.reserved:) bit, or any set bitmap bit without a matching CPUID component descriptor raises INVALID\_CONTROL\_IMAGE.

The complete \texttt{SAVE\_AREA\_SIZE\_BYTES} range reported by CPUID, every selected GS image, and every selected extension component are validated before (:term:base.terms.architectural_state:) is committed. Any address, access, or validation fault leaves the architectural register state and save area unchanged.

RESTORE treats its address-role EA as the base of the save area without first reading an EA-sized operand. It restores R0 through R15 and FLAGS from the fixed base block. A set \texttt{GS\_VALID} bit selects the corresponding GS image for validation and restoration; a clear bit preserves the current GSn. User-mode RESTORE ignores saved STATUS and DFA. Supervisor RESTORE validates and restores the complete saved supervisor STATUS, including EDEPTH and UO, together with current DFA as an architectural context transition. Ordinary WRSTATUS retains its defined (:term:base.terms.field:) restrictions. A set state-block bitmap bit restores its CPUID-described extension component. A clear bit for a described component installs that component\textquotesingle{}s initial state and makes the component clean; a restored non-initial image makes the component modified.

Before reading any extension component or changing (:term:base.terms.architectural_state:), RESTORE validates the complete base (:term:base.terms.header:) and bitmap. An unrecognized format, a nonzero (:term:base.terms.reserved:) bit, or any set bitmap bit without a matching CPUID component descriptor raises INVALID\_CONTROL\_IMAGE.

The complete \texttt{SAVE\_AREA\_SIZE\_BYTES} range reported by CPUID, every selected GS image, and every selected extension component are validated before (:term:base.terms.architectural_state:) is committed. Any address, access, or validation fault leaves the architectural register state and save area unchanged.
